\documentclass[11pt]{article}
\usepackage{amsmath,amssymb,amsthm}
\usepackage{geometry}
\geometry{margin=1in}

\newtheorem{theorem}{Theorem}
\newtheorem{lemma}{Lemma}
\newtheorem{proposition}{Proposition}
\newtheorem{corollary}{Corollary}
\theoremstyle{definition}
\newtheorem{definition}{Definition}
\newtheorem{remark}{Remark}

\newcommand{\R}{\mathbb{R}}
\newcommand{\Z}{\mathbb{Z}}
\newcommand{\Sg}{\Sigma}
\newcommand{\bd}{\partial}
\newcommand{\norm}[1]{\left\lVert #1 \right\rVert}
\newcommand{\abs}[1]{\left\lvert #1 \right\rvert}
\newcommand{\eps}{\varepsilon}
\DeclareMathOperator{\image}{im}

\begin{document}

\title{On the infimum of realized numbers for clean triangulations of the strip}
\author{}
\date{}
\maketitle

\noindent\textbf{Status: Partial.}
The claimed answer is \emph{correct}: the infimum of the set of realized numbers
is $\sqrt 3$.  This document proves rigorously and completely the two
ingredients
\begin{itemize}
\item[(U)] \emph{(Upper bound, $\inf\le\sqrt3$.)} Every $\beta>\sqrt3$ is realized.
This is proved by an explicit construction, fully verified.
\item[(L$'$)] \emph{(A clean nonsharp lower bound, $\inf>0$.)} No $\beta$ with
$\beta\le 1$ is realized; more generally a quantitative obstruction is proved.
\end{itemize}
The \emph{sharp} lower bound (L) ``no $\beta<\sqrt3$ is realized'', which together
with (U) yields $\inf=\sqrt3$, is the content of a deep theorem (the resolution of
the Halpern--Weaver conjecture). Section~\ref{sec:lower} reduces (L) to a single
clean geometric statement about embedded squeezed M\"obius bands (Lemma~\ref{lem:key}),
gives a complete proof of the weaker bound (L$'$), and explains precisely which
step in the sharp argument is the deep one. The construction in (U) is the optimal
``triangular M\"obius band'', and it is shown that the constructed bands converge
to the doubled equilateral triangle as $\beta\downarrow\sqrt3$, which is why the
infimum cannot be lowered.

\bigskip

\section{Reduction to embedded squeezed M\"obius bands}\label{sec:reduce}

Throughout, $\Sg=\R\times[0,1]$, $\bd\Sg=\R\times\{0,1\}$, and for $\beta>0$
\[
G_\beta(x,y)=(x+\beta,\,1-y).
\]

\subsection{The group generated by $G_\beta$ and the quotient}

\begin{lemma}\label{lem:group}
$G_\beta$ is a homeomorphism of $\Sg$ with $G_\beta^{2}(x,y)=(x+2\beta,y)$, and
$G_\beta^{k}(x,y)=(x+k\beta,\,y)$ if $k$ is even, $(x+k\beta,\,1-y)$ if $k$ is odd.
The cyclic group $\Gamma_\beta:=\langle G_\beta\rangle\cong\Z$ acts freely and
properly discontinuously on $\Sg$, preserving $\bd\Sg$. The quotient
$M_\beta:=\Sg/\Gamma_\beta$ is a M\"obius band (a connected, non-orientable
surface with one boundary circle), and the quotient projection
$\pi:\Sg\to M_\beta$ is a covering map.
\end{lemma}

\begin{proof}
$G_\beta$ is affine and invertible with inverse $G_\beta^{-1}(x,y)=(x-\beta,1-y)$,
hence a homeomorphism. Direct composition gives $G_\beta^2(x,y)
=G_\beta(x+\beta,1-y)=(x+2\beta,\,1-(1-y))=(x+2\beta,y)$, and the stated formula
for general $k$ follows by induction since composing with $G_\beta$ adds $\beta$ to
the first coordinate and toggles $y\mapsto 1-y$.

\emph{Freeness.} If $G_\beta^{k}(x,y)=(x,y)$ then the first coordinate forces
$x+k\beta=x$, i.e. $k\beta=0$; since $\beta>0$ this gives $k=0$.

\emph{Proper discontinuity.} For $(x,y)\in\Sg$ let $U$ be the open set
$\{(x',y'):\abs{x'-x}<\beta/2\}\cap\Sg$. If $G_\beta^{k}(U)\cap U\neq\emptyset$ for
some $k\ne0$ then some point has both its own first coordinate and its image's first
coordinate within $\beta/2$ of $x$; but the images' first coordinates are shifted by
$k\beta$ with $\abs{k\beta}\ge\beta$, so the two $x$-intervals of width $\beta$
centered at $x$ and $x+k\beta$ are disjoint, a contradiction. Hence $G_\beta^k(U)\cap
U=\emptyset$ for all $k\ne0$; the action is properly discontinuous, and $\pi$ is a
covering map of a manifold with boundary.

\emph{The quotient is a M\"obius band.} The set
$F=[0,\beta]\times[0,1]$ is a fundamental domain: every $\Gamma_\beta$-orbit meets
$F$ (reduce the first coordinate mod $\beta$), and two interior points of $F$ are
never identified. On $\bd F$, the identification induced by $\Gamma_\beta$ glues the
left edge $\{0\}\times[0,1]$ to the right edge $\{\beta\}\times[0,1]$ \emph{via}
$G_\beta$, i.e.
\[
(0,y)\ \sim\ G_\beta(0,y)=(\beta,\,1-y),\qquad y\in[0,1].
\]
Thus $M_\beta$ is the rectangle $[0,\beta]\times[0,1]$ with its two short ends
glued by an orientation-reversing flip $y\mapsto1-y$. This is the standard
description of a M\"obius band: it is connected, compact, has a single boundary
circle (the images of $\R\times\{0\}$ and $\R\times\{1\}$ join into one circle
because the flip exchanges the two boundary lines), and it is non-orientable
(the gluing reverses orientation). The number $\beta$ is, by construction, the
\emph{aspect ratio} of $M_\beta$: it is built from a $\beta\times1$ rectangle glued
with a flip.
\end{proof}

\subsection{Realizations are embeddings of $M_\beta$}

\begin{lemma}\label{lem:embed}
Let $f:\Sg\to\R^3$ be continuous and satisfy property (2) of the definition of
``realized'', i.e. $f(p)=f(q)\iff p=G_\beta^k(q)$ for some $k\in\Z$. Then $f$
factors as $f=\bar f\circ\pi$ where $\bar f:M_\beta\to\R^3$ is a continuous
injective map. If in addition $f$ satisfies property (1) (it is a squeeze map on a
$G_\beta$-invariant clean triangulation $\mathcal T$), then $\mathcal T$ descends
to a finite triangulation $\bar{\mathcal T}$ of $M_\beta$ and $\bar f$ is affine on
each triangle of $\bar{\mathcal T}$, decreasing the boundary edge and increasing the
two interior edges of each triangle.
\end{lemma}

\begin{proof}
For the forward direction of (2), $f$ is constant on each $\Gamma_\beta$-orbit, so
it factors through the orbit space: there is a unique map $\bar f:M_\beta\to\R^3$
with $f=\bar f\circ\pi$, and $\bar f$ is continuous because $\pi$ is a quotient map.
For the reverse direction of (2): if $\bar f(\pi p)=\bar f(\pi q)$ then
$f(p)=f(q)$, so $p=G_\beta^k(q)$ for some $k$, i.e. $\pi p=\pi q$. Hence $\bar f$ is
injective.

For property (1): since $\mathcal T$ is $G_\beta$-invariant and locally finite, and
$\Gamma_\beta$ acts properly discontinuously, the projection $\pi$ carries
$\mathcal T$ to a locally finite collection $\bar{\mathcal T}$ of triangles in the
compact surface $M_\beta$, hence a finite triangulation. (Two triangles of
$\mathcal T$ project to the same triangle of $\bar{\mathcal T}$ iff they differ by an
element of $\Gamma_\beta$.) Because $\pi$ is a local affine isometry on each triangle
and $f$ is affine on each triangle of $\mathcal T$, $\bar f$ is affine on each
triangle of $\bar{\mathcal T}$. The boundary edge of a clean triangle (the one in
$\bd\Sg$) projects to the boundary circle of $M_\beta$; the squeeze condition,
being a statement about edge lengths of a single triangle, descends verbatim:
$\bar f$ shortens the boundary edge and lengthens the two interior edges of each
triangle of $\bar{\mathcal T}$.
\end{proof}

\begin{definition}\label{def:smb}
A \emph{squeezed M\"obius band of aspect ratio $\beta$} is a continuous injective
map $\bar f:M_\beta\to\R^3$ which is affine on each triangle of some finite
triangulation $\bar{\mathcal T}$ of $M_\beta$ by ``clean'' triangles (each having
exactly one edge on $\bd M_\beta$), and which on each such triangle shortens the
boundary edge and lengthens the two interior edges.
\end{definition}

Combining Lemmas~\ref{lem:group}--\ref{lem:embed} with the converse (any squeezed
M\"obius band lifts to a realization via $f=\bar f\circ\pi$, using the
$G_\beta$-invariant lift of $\bar{\mathcal T}$) gives:

\begin{proposition}\label{prop:equiv}
$\beta>0$ is realized $\iff$ there exists a squeezed M\"obius band of aspect ratio
$\beta$. Consequently the set of realized numbers equals the set of aspect ratios of
embedded squeezed M\"obius bands.
\end{proposition}

\begin{proof}
($\Rightarrow$) Lemma~\ref{lem:embed}. ($\Leftarrow$) Given $\bar f$ and
$\bar{\mathcal T}$, lift $\bar{\mathcal T}$ along $\pi$ to a $G_\beta$-invariant
clean triangulation $\mathcal T$ of $\Sg$ (the preimage of each triangle is a
$\Gamma_\beta$-orbit of triangles; this is countable and locally finite because
$\pi$ is a covering and $\bar{\mathcal T}$ is finite). Set $f=\bar f\circ\pi$. Then
$f$ is continuous, affine and squeezing on each triangle of $\mathcal T$ (property
(1)), and $f(p)=f(q)\iff\pi p=\pi q\iff p=G_\beta^k(q)$ (property (2), using
injectivity of $\bar f$ and the definition of $\pi$). So $\beta$ is realized.
\end{proof}

Thus the problem is exactly: \emph{determine the set of aspect ratios admitting an
embedded squeezed M\"obius band, and show its infimum is $\sqrt3$.} We treat the
upper and lower bounds separately.

\section{Upper bound: every $\beta>\sqrt3$ is realized}\label{sec:upper}

We construct, for each $\beta>\sqrt3$, an explicit squeezed M\"obius band of aspect
ratio $\beta$. By Proposition~\ref{prop:equiv} this shows every such $\beta$ is
realized, hence $\inf\{\text{realized}\}\le\sqrt3$.

\subsection{The model: a doubled equilateral triangle}\label{sec:model}

Fix the equilateral triangle $T_0\subset\R^2\times\{0\}\subset\R^3$ with vertices
\[
P_1=(0,0,0),\qquad P_2=(1,0,0),\qquad P_3=\Big(\tfrac12,\tfrac{\sqrt3}{2},0\Big),
\]
so $T_0$ has unit side length and centroid $O=(\tfrac12,\tfrac{1}{2\sqrt3},0)$.
The ``triangular M\"obius band'' is, informally, the boundary strip obtained by
running a thin ribbon three times around the edges of $T_0$, the ribbon making a
$60^\circ$ turn at each corner, so that after three corners ($3\times60^\circ
=180^\circ$ of accumulated turning of the ribbon's centerline relative to its
transverse direction) the ribbon closes up with a half-twist. We make this precise
and verify it is an \emph{embedding} with the squeeze property, for aspect ratios
arbitrarily close to (and above) $\sqrt3$.

\subsection{Construction}\label{sec:construct}

Fix $\beta>\sqrt3$ and a small parameter $\delta>0$ to be chosen
($\delta\to0$ as $\beta\downarrow\sqrt3$). We work with the fundamental rectangle
\[
R=[0,\beta]\times[0,1],\qquad (0,y)\sim(\beta,1-y),
\]
which by Lemma~\ref{lem:group} is $M_\beta$.

\paragraph{Centerline polygon.}
Let $h>0$ and consider the planar isoceles ``arrowhead'' obtained by pushing $T_0$
to a degenerate doubled triangle: define the three \emph{rails}
\[
L_1,\,L_2,\,L_3
\]
to be the three sides of an equilateral triangle $T_\delta$ of side $\ell=\ell(\delta)$
lying in the horizontal plane $z=0$, and \emph{lift} the ribbon a height
$\pm\delta$ alternately so that the surface is embedded rather than flat. Concretely,
subdivide the boundary parameter $x\in[0,\beta]$ into three equal arcs of length
$\beta/3$ by the points
\[
x_0=0<x_1=\tfrac\beta3<x_2=\tfrac{2\beta}3<x_3=\beta .
\]
On the $j$-th arc $[x_{j-1},x_j]$ ($j=1,2,3$) the band will run alongside the side
$L_j$ of $T_\delta$.

\paragraph{Placement map on the centerline.}
Let $\gamma:[0,\beta]\to\R^3$ be the closed (after the flip-gluing) piecewise-linear
centerline that traverses the three sides $L_1,L_2,L_3$ of $T_\delta$ in order, at
heights alternating in a small band: explicitly, writing the corners of $T_\delta$
as $Q_1,Q_2,Q_3$ (with $Q_{j}$ shared by $L_{j-1}$ and $L_j$, indices mod $3$),
\[
\gamma(x)=\Big(1-s_j(x)\Big)\,A_j+s_j(x)\,B_j,\qquad x\in[x_{j-1},x_j],
\]
where $s_j(x)=\frac{x-x_{j-1}}{x_j-x_{j-1}}\in[0,1]$ parametrizes the $j$-th side
linearly from its initial corner $A_j=Q_j$ to its terminal corner $B_j=Q_{j+1}$.
This $\gamma$ is the centerline of an equilateral triangle of side $\ell$, traversed
once; its total length is $3\ell$.

\paragraph{Ribbon (the map $\bar f$).}
At each point of the $j$-th side, attach a transverse unit interval; to obtain a
\emph{M\"obius} band the transverse direction must rotate by a half turn over the
whole loop. We realize this by letting the transverse direction at side $L_j$ be a
fixed unit vector $u_j\in\R^3$, chosen so that consecutive $u_j$ differ by a
$60^\circ$ rotation \emph{in the plane perpendicular to the centerline}, and so that
$u_{3}$ continues to $-u_1$ across the gluing corner (the half-twist). Define, for
$(x,y)\in R$,
\begin{equation}\label{eq:ribbon}
\bar f(x,y)=\gamma(x)+\big(y-\tfrac12\big)\,w_j(x),\qquad x\in[x_{j-1},x_j],
\end{equation}
where $w_j(x)$ is the unit transverse field on the $j$-th panel, defined below so
that \eqref{eq:ribbon} is continuous across the corners and across the gluing.

\paragraph{Choice of transverse fields and the half-twist.}
Place $T_\delta$ so that its plane is $z=0$ and its centroid is the origin. Let
$n=(0,0,1)$ be the upward normal and let $t_j$ be the unit tangent along side $L_j$
(pointing in the direction of increasing $x$). On panel $j$ we set
\begin{equation}\label{eq:wfield}
w_j(x)=\cos\!\big(\theta_j(x)\big)\,n+\sin\!\big(\theta_j(x)\big)\,t_j,
\end{equation}
a unit vector in the plane spanned by the normal $n$ and the tangent $t_j$, where
$\theta_j(x)$ increases monotonically from $\theta_j(x_{j-1})$ to
$\theta_j(x_j)$ so that the total increment of $\theta$ around the loop is exactly
$\pi$:
\[
\theta_1(0)=0,\qquad \theta_3(\beta)=\pi,\qquad
\theta_j(x_j)=\theta_{j+1}(x_j)\ (\text{continuity at corners}),
\]
with each panel contributing $\pi/3$:
$\theta_j(x_j)-\theta_j(x_{j-1})=\pi/3$. Concretely
\[
\theta_j(x)=\frac{\pi}{3}\,(j-1)+\frac{\pi}{3}\,s_j(x),\qquad x\in[x_{j-1},x_j].
\]
Thus $w$ starts pointing straight up ($\theta=0$) and ends pointing straight down
relative to its starting frame ($\theta=\pi$, i.e. $\cos\pi=-1$); this is exactly the
half-twist that makes \eqref{eq:ribbon} descend to the M\"obius band $M_\beta$ rather
than to a cylinder.

\paragraph{Closing up (the gluing).}
At the gluing edge we must have
$\bar f(0,y)=\bar f(\beta,1-y)$. From \eqref{eq:ribbon},
$\bar f(0,y)=\gamma(0)+(y-\tfrac12)w_1(0)$ with $w_1(0)=n$ (as $\theta=0$), and
$\bar f(\beta,1-y)=\gamma(\beta)+(\tfrac12-y)w_3(\beta)$ with
$w_3(\beta)=\cos(\pi)n+\sin(\pi)t_3=-n$ (as $\theta=\pi$). Since $\gamma$ is a closed
triangle, $\gamma(0)=\gamma(\beta)=Q_1$. Hence
\[
\bar f(0,y)=Q_1+(y-\tfrac12)n,\qquad
\bar f(\beta,1-y)=Q_1+(\tfrac12-y)(-n)=Q_1+(y-\tfrac12)n,
\]
so the two agree: the gluing $(0,y)\sim(\beta,1-y)$ is respected, and $\bar f$
descends to a continuous map $M_\beta\to\R^3$.

\paragraph{Clean triangulation and squeeze.}
Triangulate $R$ as follows. On panel $j$, subdivide $[x_{j-1},x_j]\times[0,1]$ into
$N$ vertical strips of width $\tfrac{\beta}{3N}$ and triangulate each unit square
strip by its two diagonals so that every triangle has \emph{exactly one} edge on
$\bd R=\R\times\{0,1\}$ (the standard ``alternating up/down'' triangulation of a thin
rectangle: bottom-base triangles have their base on $y=0$ and apex on $y=1$,
top-base triangles vice versa). Each such triangle is a clean triangle, and the
collection is $G_\beta$-invariant once lifted (it is invariant under the gluing by
construction because the subdivision points are placed symmetrically with respect to
the flip $y\mapsto 1-y$, which we may arrange by choosing $N$ subdivision points
symmetric under $y\mapsto1-y$). The map $\bar f$ in \eqref{eq:ribbon} is not affine on
these triangles (because $w_j(x)$ varies with $x$), so we replace $\bar f$ by its
\emph{affine interpolant} $\bar f_N$: $\bar f_N$ agrees with $\bar f$ at all vertices
of the triangulation and is affine on each triangle. As $N\to\infty$, $\bar f_N\to
\bar f$ uniformly with derivatives converging, by piecewise-linear approximation of
the smooth-on-each-panel map $\bar f$.
\end{document}
